\section{Demo trực tiếp}\label{sec:demo}
Kể demo như \key{câu chuyện của một vị khách}: chị ghé Họa., tìm tranh bằng lời và bằng giọng nói, đưa ảnh một bức tranh rồi ảnh phòng khách, đặt mua; người quản lý chuyển đơn đi; chị hỏi trợ lý và tò mò thử xem đơn của người khác; cuối cùng ta rút phích AI. Năm flow, cùng thứ tự với mục F của showcase, tổng khoảng \key{6 phút}. Trong kịch bản 7 bước (bước 6) chỉ dùng một phần Flow 1, 3 và 4; các flow còn lại dùng khi thầy muốn xem thêm. Mọi câu tìm đều bằng \key{tiếng Anh}.

\begin{note}[warn]{Ba điều cần biết trước khi demo}
\begin{itemize}
\item Gemma có thể trả lời \key{hơi khác giữa các lần} (câu lý do, thứ tự các bức phía sau). Phần ``Mong đợi'' chỉ là tinh thần, đừng đọc nguyên văn.
\item Giỏ của C001 có sẵn \key{một bức} (Woman with a Parasol, 89 USD), nên tổng tiền lớn hơn món vừa thêm. Cứ nói ``trong giỏ có sẵn một bức''.
\item Muốn cho thầy thấy AI có tham gia hay không: \key{Shift+D} bật chế độ nhà phát triển, hiện nhãn ``GEMMA 4'' hoặc ``KEYWORD FALLBACK'' và thời gian. Bật khi cần, rồi tắt.
\end{itemize}
\end{note}

\subsection{Flow 1 · Tìm bằng lời và bằng giọng nói \normalfont\small(≈ 1 phút; UC02, UC03)}
\lab{Thao tác}
Tab cửa hàng, bấm \emph{Explore the collection}. Gõ \code{blue landscape paintings}, Enter. Rồi bấm micro trong ô tìm, nói rõ \code{Monet paintings under 100 dollars}, bấm dừng; chữ hiện trong ô \emph{We heard}, sửa nếu sai, bấm \emph{Search}.
\lab{Nói}
Một vị khách ghé Họa. Chị không nhớ tên bức tranh, chỉ nhớ đó là phong cảnh, màu xanh. Chị gõ đúng như vậy, và mỗi bức hiện ra kèm một dòng lý do vì sao được chọn. \cue{bấm micro, nói tiếng Anh} Lần này chị nói: tranh của Monet, dưới một trăm đô. Giọng nói được đổi thành chữ ngay trên chiếc laptop này, âm thanh không đi đâu cả; chữ hiện ra để chị sửa nếu máy nghe sai. Ba bức đầu đúng là ba bức Monet dưới một trăm đô.
\lab{Mong đợi}
Chữ: khoảng mười hai bức, bức đầu thường là Woman with a Parasol, mỗi thẻ có dòng lý do. Giọng nói: Impression, Sunrise (79 USD), Water Lilies (69), Woman with a Parasol (89), rồi vài bức Ấn tượng khác dưới 100 USD (cửa hàng chỉ có ba bức Monet đủ điều kiện).
\lab{Nếu lỗi}
Nhãn ``KEYWORD FALLBACK'': Gemma không trả lời (mục~\ref{sec:khonggemma}). \emph{Voice search is unavailable}: Whisper chưa chạy; gõ câu thay vì nói (mục~\ref{sec:suco}). Máy nghe ``Monet'' thành ``Money'': sửa trong ô \emph{We heard}, và nói luôn: đây là lý do khách được xem và sửa chữ.

\subsection{Flow 2 · Đưa ảnh, rồi treo tranh lên tường \normalfont\small(≈ 1 phút 30; UC04, UC14)}
\lab{Thao tác}
Ở trang Collection bấm \emph{find paintings similar to a photo}, chọn \file{Q02.jpg} trên Desktop. Rồi kéo thả \file{room1.jpg} vào trang (hoặc Cmd+V): trang chuyển sang \emph{Find for my room}, bức hợp nhất đã treo sẵn; bấm một bức khác trong dải \emph{Picked for this wall}.
\lab{Nói}
Chị có một tấm ảnh chụp bức tranh ở bảo tàng, đây là bức Sóng lừng. Gemma nhìn ảnh và tả nó bằng ba nhãn: phong cách, chủ đề, màu; cửa hàng tìm những bức có nhãn giống nhất, và đúng bức đó đứng đầu. \cue{kéo ảnh phòng vào} Rồi chị muốn biết tranh có hợp phòng khách không. Chị đưa ảnh căn phòng: bức hợp nhất được treo ngay lên tường trong ảnh, kèm một câu giải thích nhắc tới những thứ có thật trong phòng. \cue{bấm bức khác} Bấm bức khác thì nó treo lên ngay. Ảnh phòng chỉ nằm trong trình duyệt, không được lưu lên máy chủ.
\lab{Mong đợi}
Ảnh Q02: The Great Wave off Kanagawa đứng đầu. Phòng: ``Looking at your room\ldots'' vài giây, rồi tranh trên tường và câu \emph{Why it fits}.
\lab{Nếu lỗi}
Chờ tối đa 15 giây; quá hạn thì cửa hàng tự chọn theo màu và dựng câu giải thích từ thông tin của tranh. Ảnh lỗi định dạng: dùng một trong ba phòng mẫu ở cột bên phải.

\subsection{Flow 3 · Mua, xem danh mục, rồi giao hàng \normalfont\small(≈ 1 phút 30; UC05, UC07, UC12, UC13)}
\lab{Thao tác}
\begin{enumerate}[leftmargin=1.5em]
\item Mở một bức (ví dụ The Great Wave). \key{Chỉ vào dòng nhỏ trên tên tranh}: \emph{Fine-art prints ▪ Ukiyo-e ▪ 1831}, tức tên danh mục ▪ phong cách ▪ năm. Một cái khung thì dòng này là \emph{Frames ▪ vật liệu}.
\item Chọn một khung, \emph{Add to cart}, mở \emph{Cart}. Chọn địa chỉ, \emph{Cash on delivery}, \emph{Place order}.
\item Bấm tên ở góc phải trên, chọn \emph{Gallery staff}; vào \emph{Back office}. Tab \emph{Paintings \& frames}: dưới tên mỗi sản phẩm có tên danh mục (``Fine-art prints ▪ họa sĩ'' hoặc ``Frames''). Kéo xuống form \emph{Add a painting}: ô \emph{Category} chỉ liệt kê danh mục chứa tranh.
\item Tab \emph{Orders}: ở đơn vừa đặt bấm \emph{Mark processing}. Rồi đổi lại về C001.
\end{enumerate}
\lab{Nói}
Chị chọn bức này. Dòng chữ nhỏ trên tên tranh là danh mục của nó, chính là dòng trong bảng danh mục em vừa cho thầy xem. \cue{thêm vào giỏ, đặt hàng} Chị thêm tranh và khung vào giỏ rồi đặt hàng. Phía sau, cửa hàng đọc lại giá và kho, tạo đơn với giá đã chốt, trừ kho rồi làm rỗng giỏ; tổng tiền đã trừ mười phần trăm cho khách VIP và cộng tám đô phí giao. \cue{đổi sang Gallery staff} Bây giờ em là người quản lý. Bảng sản phẩm cũng ghi danh mục của từng món, và khi thêm tranh mới, em chỉ chọn được danh mục dành cho tranh, nên không thể xếp nhầm một bức tranh vào mục khung. \cue{Mark processing} Đơn chỉ đi từng bước: mới đặt, đang xử lý, đã gửi, đã giao; mỗi đơn chỉ có đúng một nút cho bước kế tiếp.
\lab{Mong đợi}
``ORDER PLACED: Your order \#\ldots{} is confirmed''. Tổng $=$ tạm tính $-$ 10\% $+$ 8 USD. Sau khi bấm, đơn chuyển sang \emph{Processing} và nút đổi thành \emph{Mark shipped}.
\lab{Nếu lỗi}
Báo hết hàng: dữ liệu còn dư từ lần tập trước; đặt lại dữ liệu mẫu (mục~\ref{sec:reset}).

\subsection{Flow 4 · Hỏi trợ lý, rồi thử xem trộm \normalfont\small(≈ 1 phút; UC10, UC09, UC11)}
\lab{Thao tác}
Là C001: bấm nút chat tròn ở góc phải dưới, hỏi \code{where is my latest order?}. Rồi mở \emph{My orders}; ở ô \emph{Find an order} gõ \code{20261003} (đơn của C002), bấm \emph{Look up}. Cuối cùng bấm \emph{Cancel order} ở đơn vừa đặt.
\lab{Nói}
Chị hỏi trợ lý: đơn mới nhất của tôi đâu rồi? Trợ lý tự gọi công cụ tra đơn và trả lời về đúng đơn vừa đặt, đang được xử lý. Cửa hàng luôn tra theo tài khoản đang đăng nhập, không theo những gì khách gõ. \cue{tra 20261003} Giờ chị tò mò gõ mã đơn của một người khác. Cửa hàng trả lời: không tìm thấy trong tài khoản của bạn, giống hệt như khi mã đó không tồn tại, nên không ai dò ra được mã đơn nào là thật. \cue{Cancel order} Chị đổi ý và hủy đơn vừa đặt: được, vì đơn chưa gửi đi. Đơn đã gửi thì không còn nút hủy.
\lab{Mong đợi}
Trợ lý nói về đơn vừa đặt (đang xử lý). Tra 20261003: \emph{``We can't show order \#20261003 \ldots{} was not found in your account.''} Hủy: đơn chuyển sang \emph{Cancelled}; đơn 20261001 (đã gửi) không có nút hủy.

\subsection{Flow 5 · Rút phích Gemma \normalfont\small(≈ 1 phút; NFR02, NFR12)}
\lab{Thao tác}
Sang tab bản tắt Gemma (\code{http://127.0.0.1:8687/?debug=1}). Gõ lại \code{blue landscape paintings}; tải \file{Q02.jpg} ở \emph{find paintings similar to a photo}; hỏi trợ lý \code{find my order 20261001}.
\lab{Nói}
Câu hỏi thầy có thể đang nghĩ: nếu AI hỏng thì sao? Đây là một bản của cửa hàng mà đường tới Gemma đã bị cắt. Tìm bằng chữ vẫn ra, chỉ khác nhãn: lần này luật từ khóa làm việc, và còn nhanh hơn. Với ảnh thì khác: cửa hàng chỉ còn đọc được màu trung bình nên trả về bức sai, đúng như em đã đo, Gemma giúp rõ nhất ở ảnh. Trợ lý vẫn tra được đơn nhờ luật. AI hỏng, cửa hàng vẫn bán được hàng.
\lab{Mong đợi}
Chữ: vẫn ra kết quả, nhãn ``KEYWORD FALLBACK'' (bản này mở sẵn chế độ nhà phát triển). Ảnh: kết quả xếp theo màu, không phải The Great Wave. Trợ lý: vẫn trả lời đúng đơn; dưới câu trả lời ghi ``Fallback ▪ no AI'' và ``Tool ▪ find\_order''.

\subsection{Khi Gemma không chạy}\label{sec:khonggemma}
Nói câu mở ở mục~\ref{sec:gemma} trước khi bắt đầu demo, rồi đi các flow như sau. Flow 5 bỏ, vì cửa hàng chính lúc này đã là bản dự phòng.
\begin{center}\small
\begin{tabular}{@{}P{2cm}P{6.6cm}P{7.2cm}@{}}
\toprule
\textbf{Flow} & \textbf{Chuyện gì xảy ra} & \textbf{Làm và nói gì}\\
\midrule
1 · chữ, giọng & Vẫn chạy: luật từ khóa hiểu câu; Whisper vẫn chạy trên laptop & Bật Shift+D, chỉ nhãn ``KEYWORD FALLBACK'': ``đây là đường lui em thiết kế sẵn''. Câu Monet vẫn ra đúng ba bức Monet dưới 100 USD.\\
2 · ảnh, phòng & Ảnh chỉ còn đọc màu, kết quả kém; phòng xếp theo màu, lý do dựng từ thông tin của tranh & Làm nhanh một lần để thấy kết quả kém, rồi mở showcase F (``Show a photo, then hang it on your wall'') hoặc \file{figures/demo/search_image.png}, \file{room_preview.png}: ảnh chụp khi có Gemma.\\
3 · mua, giao & Không ảnh hưởng (NFR12) & Làm đầy đủ: đây là phần chứng minh cửa hàng vẫn bán được.\\
4 · trợ lý & Trợ lý trả lời bằng luật (chế độ nhà phát triển ghi ``Fallback ▪ no AI''); quyền xem đơn y nguyên & Hỏi \code{find my order 20261001} hoặc \code{where is my latest order?}, vì luật chỉ nhận mã đơn và vài từ khóa; phần tra đơn của người khác vẫn làm như cũ.\\
\bottomrule
\end{tabular}
\end{center}
Nếu được hỏi ``vậy số đo có Gemma ở đâu?'': mục G của showcase và bài §3.15 (đo ngày 30/10).

\subsection{Nếu thầy muốn xem VP sinh cơ sở dữ liệu ngay tại chỗ}\label{sec:vpgen}
Bài chưa có ảnh chụp trình hướng dẫn \emph{Generate Database} của VP (mục này trong bảng nghiệm thu còn ``một phần''); file \file{demo/db/vp_generated_ddl.sql} là do bộ sinh SQL của chính VP viết ra khi dựng project. Nếu thầy muốn tận mắt thấy, làm như sau (khoảng 2 phút). Nói thật trước: \say{phần này em làm trực tiếp, bài chưa có ảnh chụp từng bước.}
\begin{enumerate}[leftmargin=1.5em]
\item Giữ lại bản cũ: \code{cp demo/db/vp_generated_ddl.sql ~/Desktop/vp_generated_ddl.bak.sql}.
\item Mở \file{vp/tieuluan_design.vpp} bằng \key{bản dùng thử} (\file{Visual Paradigm.app}), mở sơ đồ \vp{4.2 Database - entity relationship diagram}.
\item Menu \emph{Tools $\to$ DB $\to$ Generate Database}. Giữ cơ sở dữ liệu là \key{SQLite}, chọn \emph{Generate DDL} ra file, lưu đè lên \file{demo/db/vp_generated_ddl.sql} (đó là file công cụ so sánh đọc).
\item Ở thư mục dự án chạy \code{python3 tools/compare_schema.py}. Mong đợi: dòng ``Same tables, columns, types, NOT NULL, primary keys and UNIQUE: True (\factCheckTables{} tables, \factCmpColumns{} columns compared)'' và câu ghi rằng VP không viết khóa ngoại và CHECK.
\item Đóng VP, chọn \emph{Don't Save}.
\end{enumerate}
\lab{Nói}
Em bấm sinh cơ sở dữ liệu ngay trong Visual Paradigm, từ chính sơ đồ quan hệ thực thể. Rồi em cho công cụ so file vừa sinh với cơ sở dữ liệu cửa hàng đang chạy: giống nhau về bảng, cột, kiểu dữ liệu, khóa chính và ràng buộc duy nhất. Chỉ khác là bộ sinh của VP cho SQLite không viết khóa ngoại và ràng buộc kiểm tra, nên công cụ của em thêm chúng vào từ cùng mô hình.
\lab{Nếu lỗi}
Có dòng \code{DIFFERENCE}: đọc to dòng đó và giải thích; muốn quay lại bằng chứng cũ thì \code{cp ~/Desktop/vp_generated_ddl.bak.sql demo/db/vp_generated_ddl.sql} rồi chạy lại lệnh so. Bản dùng thử không mở được: mở \file{build/schema_comparison.txt} và \file{demo/db/vp_generated_ddl.sql} để thầy xem kết quả lần trước.
